\section{The Proposed Method: \method{}}
\label{sec:method}

% Recommended structure (adapt to fit your contribution):
%   III.A  Overview (with Figure 1)
%   III.B  Formal problem statement / preliminaries
%   III.C  Core algorithm or framework with pseudocode
%   III.D  (Optional) Complexity / theoretical analysis
%   III.E  Implementation details (parameters, defaults)
%
% Re-section as appropriate. Not every paper needs all of these
% subsections; not every paper has a multi-stage pipeline.

\subsection{Overview}
\label{sec:method:overview}

Figure~\ref{fig:overview} shows the high-level structure of \method{}.
\textbf{[REPLACE with a 2--3 sentence summary of the framework and where
the novelty lies. Adapt the figure below to match.]}

\begin{figure}[!t]
    \centering
    % Placeholder TikZ overview diagram.
    % Replace with your real figure (TikZ preferred over raster).
    \begin{tikzpicture}[
        node distance=0.6cm,
        block/.style={draw, rectangle, rounded corners=2pt,
                      minimum width=2.2cm, minimum height=0.7cm,
                      align=center, font=\small},
        arrow/.style={-Stealth, thick}
    ]
        \node[block, fill=blue!10]   (in)   {Input\\(\eg{} RTL or netlist)};
        \node[block, right=of in, fill=green!10]  (s1) {Stage 1};
        \node[block, right=of s1, fill=orange!10] (s2) {Stage 2\\(novel)};
        \node[block, right=of s2, fill=red!10]    (s3) {Stage 3};
        \node[block, right=of s3]                  (out){Output\\(\eg{} placed design)};
        \draw[arrow] (in)--(s1);
        \draw[arrow] (s1)--(s2);
        \draw[arrow] (s2)--(s3);
        \draw[arrow] (s3)--(out);
    \end{tikzpicture}
    \caption{Overview of \method{}. Stages~1 and~3 follow standard EDA
    practice; Stage~2 (highlighted) is the novel component proposed in
    this paper, described in Section~\ref{sec:method:core}.}
    \label{fig:overview}
\end{figure}

\subsection{Preliminaries and Notation}
\label{sec:method:prelim}

\textbf{[REPLACE]} We adopt the following notation throughout the paper.
Let \(D\) denote a [design / netlist / placement instance] with [N] cells
and [M] nets. The objective is to minimize
\(\Phi(D) = \alpha \cdot \mathrm{area}(D) + \beta \cdot \mathrm{delay}(D)
+ \gamma \cdot \mathrm{power}(D)\), subject to [constraints]. The
hyperparameters \(\alpha, \beta, \gamma\) are set in
Section~\ref{sec:method:impl}.

\subsection{Core Algorithm}
\label{sec:method:core}

The core of \method{} is summarized in Algorithm~\ref{alg:method}.

\begin{algorithm}[!t]
\caption{\method{}: high-level procedure.}
\label{alg:method}
\begin{algorithmic}[1]
\REQUIRE Design \(D\), parameters \(\theta = (\alpha,\beta,\gamma)\)
\ENSURE Optimized design \(D^*\)
\STATE \(D_0 \leftarrow \text{Stage1}(D)\)
\STATE \(k \leftarrow 0\); \(\Phi_{\text{prev}} \leftarrow \infty\)
\WHILE{not converged}
    \STATE \(D_{k+1} \leftarrow \text{Stage2}(D_k, \theta)\)
        \COMMENT{novel step (Sec.~\ref{sec:method:core})}
    \STATE \(\Phi_k \leftarrow \Phi(D_{k+1})\)
    \IF{\(|\Phi_{\text{prev}} - \Phi_k| < \varepsilon\)}
        \STATE \textbf{break}
    \ENDIF
    \STATE \(\Phi_{\text{prev}} \leftarrow \Phi_k\); \(k \leftarrow k+1\)
\ENDWHILE
\STATE \(D^* \leftarrow \text{Stage3}(D_{k+1})\)
\RETURN \(D^*\)
\end{algorithmic}
\end{algorithm}

\textbf{Stage 1 — [REPLACE name].} [Two-three sentences explaining the
stage. Cite related work used at this stage.]

\textbf{Stage 2 — [REPLACE name, the novel part].} The key insight is
that [the insight]. We formalize this as follows:

\begin{equation}
\label{eq:novel}
f(D) \;=\; \sum_{e \in E} w(e) \cdot \phi\!\left(\, x(e) \,\right)
       \;+\; \lambda \, R(D),
\end{equation}

\noindent where \(w(\cdot)\) is the [weight function], \(\phi(\cdot)\) is
[smoothing function], and \(R(\cdot)\) is the [regularization term]
weighted by \(\lambda\). The function \(f(D)\) is solved by [optimization
method, \eg{} L-BFGS / coordinate descent / dynamic programming], yielding
[output property].

\textbf{Stage 3 — [REPLACE name].} [Two-three sentences.]

\subsection{Complexity Analysis}
\label{sec:method:complexity}

\textbf{[REPLACE]} The dominant cost of Algorithm~\ref{alg:method} is
Stage~2, which performs \(O(\log(1/\varepsilon))\) iterations of
\(O(n \log n)\) work each, where \(n\) is the number of cells in the
design and \(\varepsilon\) is the convergence tolerance. Total time
complexity is therefore \(O(n \log n \cdot \log(1/\varepsilon))\).
Memory usage is \(O(n + m)\) where \(m\) is the number of nets.
We empirically validate this scaling in
Section~\ref{sec:experiments:scaling}.

\subsection{Implementation Details}
\label{sec:method:impl}

\textbf{[REPLACE]} \method{} is implemented in approximately [N] lines of
[Python / C++] on top of [base tool, \eg{} OpenROAD]. The default
hyperparameters are \(\alpha = 1.0\), \(\beta = 0.5\), \(\gamma = 0.1\),
and \(\varepsilon = 10^{-3}\); these were tuned on the validation subset
described in Section~\ref{sec:experiments:setup} and held fixed across
all reported experiments. The convergence tolerance \(\varepsilon\)
typically triggers in 8--12 outer iterations on designs we tested.
